\documentclass[a4paper, 10pt]{article}
\usepackage[top=1cm,bottom=1cm,left=1cm,right=1cm]{geometry}

\usepackage[ngerman,english]{babel}
\usepackage[ansinew]{inputenc}
\usepackage[colorlinks=true, urlcolor=blue]{hyperref}

\usepackage{graphicx}

\usepackage{enumitem}
\usepackage{multicol}
\usepackage{amssymb}
\usepackage{amsmath}
\usepackage{amstext}
\usepackage{amsfonts}
\usepackage{eurosym}

\begin{document}
{
	Lineare Differentialgleichung
	\[ \begin{cases} \ddot x + 2 \dot x + 17 x &= 2 \cdot \sin(5\cdot t)\,,\\ x(\pi) &= 0\,,\\ \dot x(\pi) &= 1\,.\end{cases} \]
	Konvergenz von Potenzreihen: Definitionsbereich der folgenden Funktionen?
	\begin{eqnarray*}
		\begin{array}{l r c l}
			(a) \, & f(x) &=& \sum\limits_{n=1}^\infty 2^{n^2} x^{2n+1} \\[2ex]
	(b) \, & f(x) &=& \sum\limits_{n=1}^\infty \dfrac{3^n}{\sqrt{ 2^n\left( 3n-1\right) }} (x-1)^n
	\end{array}
	\end{eqnarray*}
	Fourierreihe: Wie lauten \(a_5\) und \(b_5\)? \[ f|_T: t\mapsto |3-t|, \quad t\in [0,6)\,.\]
	Fourierreihe: Wie lautet die Fourierreihe des folgenden Signals \[ f|_T: t\mapsto \max\left( \sin(t),0\right), \quad t\in [0,2\pi)\,.\]


}
\end{document}
